\section{Trust: what exactly is being checked?}
\label{sec:trust}

\subsection{Three layers that are not interchangeable}

It is essential to keep apart the mathematical soundness of a certificate
format, the software correctness of a checker, and a proof accepted by Lean.
The propositions in this document explain why each certificate implies the
statement it claims. Tests provide evidence about the Python implementations.
Neither supplies a Lean theorem. A completed tactic must either reconstruct a
Lean proof or use a formally verified checker whose result is established
inside Lean.

\begin{proposition}[Certificate-format soundness is not implementation
correctness]
Let $\mathcal{C}$ be a certificate format with a proved soundness statement
``$\mathrm{check}(P,C)=\texttt{true}\Rightarrow P$ holds'', and let $K$ be a
program intended to compute $\mathrm{check}$. Evidence that $K$ accepts only
certificates satisfying $\mathrm{check}$ is evidence about $K$, obtained by
testing. It does not follow from the soundness statement, and the soundness
statement does not follow from it.
\end{proposition}

This is laboured deliberately, because the gap is where an automation project
most easily deceives itself. The prototypes here replay every stored
certificate with the standard library alone, and several run their checkers in
a direction algorithmically opposite to their producers --- a Bernstein checker
that \emph{expands} a claimed basis representation and compares it against an
independently computed affine transform shares no conversion formula with the
search that proposed it. That is real, and it is still not kernel verification.

\subsection{The acceptance judgement}

An accepted source theorem must have:

\begin{enumerate}[leftmargin=1.4em]
\item no unresolved expression or universe metavariables;
\item no \code{sorry} dependency anywhere in its transitive closure;
\item no unexpected new axioms relative to the approved environment;
\item every guard and side condition discharged, not merely recorded;
\item a proof of \emph{the exact original goal}, not of a restated one.
\end{enumerate}

The last point deserves emphasis because it is the one a certificate pipeline
gets wrong most easily.

\begin{invariant}[Certificate binding]
\label{inv:binding}
A certificate is checked against the problem derived from the \emph{source}
--- the same variable mapping, types, assumptions, and target. The checker
receives the original target and hypotheses externally. It must never accept a
problem statement supplied by the certificate itself. A certificate for
another polynomial, another box, or a different quantifier scope is not
interchangeable with the one required.
\end{invariant}

A correct polynomial identity about the wrong reification is not a proof of
the user's goal. A certificate must identify the actual hypotheses it used,
not merely provide polynomial data that resembles them. The prototypes test
this directly: removing a necessary constraint from the problem makes replay
fail, and a certificate re-attached to a changed target is rejected.

\subsection{Kernel-only and native-checked are different profiles}

The Lean kernel verifies derivations relative to the axioms they use, so a
strict automation policy must inspect the transitive axiom dependencies of the
final proof and of its auxiliary declarations. A practical whitelist may allow
the project's usual logical axioms --- propositional extensionality, quotient
soundness, and classical choice where permitted --- while rejecting admitted
proofs and compiler-result axioms.

The current tactic reference documents that \code{native\_decide} and
\code{bv\_decide} rely on native code generation and associated trusted
results~\cite{tactic-reference,validating-proofs}. Consequently ``the SAT
solver returns a checked certificate'' does not by itself mean ``no compiler
trust is added''. \forge{} should place that route behind an opt-in trust mode
rather than enabling it under a kernel-only label, and the active profile must
be visible in the trace and in the output metadata.

Two further details matter for anyone implementing the audit. Native
evaluation in recent Lean versions can introduce a \emph{separate generated
axiom per computation}, so a single-name blacklist such as ``reject
\code{Lean.trustCompiler}'' is not a sufficient
policy~\cite{native-infra,native-rfc}. And the audit must follow the actual
pinned Lean version rather than a remembered one; the set of axioms a given
tactic can introduce is a property of the release.

A kernel-only bit-vector route would require replaying a supported certificate
checker through a reduction mechanism that introduces no native-result axioms,
and measuring the resulting cost. That route is proposed here, not implemented.
Lean4Lean is the relevant reference point for what independent kernel-level
assurance would mean~\cite{lean4lean}.

\subsection{Diagnostic holes are not accepted proofs}

The inspected \grind{} configuration includes an option concerning admitted
branches in trace-oriented diagnostics~\cite{grind-config}. That is not a claim
that ordinary successful \grind{} proofs are unsound. It is a reminder that a
coordinator must never mistake a diagnostic artefact for a finished proof: the
root declaration's dependencies must be checked for admitted terms even when an
inner tool reports a successful diagnostic operation.

Absence of the token \code{sorry} in a source file is likewise not evidence
that the file compiles, and it is not evidence that the declarations it
contains are proved. A benchmark harness should detect an admitted declaration
in a tactic's output and record it as a rejection rather than a success.

\subsection{Soundness invariants to assert in code}

\begin{invariant}[Context validity]
Every proved fact is accompanied by a proof whose free variables and local
instances belong to its recorded context, or by an explicit abstraction and
instantiation that transports it into the consuming context.
\end{invariant}

\begin{invariant}[Candidate quarantine]
No candidate theorem, model equation, numerical conjecture, or unsolved proof
hole is admitted to the proved-fact state. Exploratory assumptions never serve
as proof evidence.
\end{invariant}

\begin{invariant}[Induction discipline]
An induction hypothesis is used only with the scope and smaller-argument
relationship supplied by its recursor. Generalised parameters may vary; the
structural induction index may not be replaced arbitrarily.
\end{invariant}

\begin{invariant}[Fail-closed arithmetic]
Numerical tolerances may guide search. Accepted equalities, coefficient signs,
partition boundaries, and proof-rule premises are checked exactly.
\end{invariant}

\begin{invariant}[Acyclic justification]
A lemma is not available while its own certificate is being checked, and a
bundle of lemmas is accepted only in a dependency order in which no entry
refers to a later one. Name shadowing of a defining equation or of the
induction hypothesis is rejected outright.
\end{invariant}

Invariant~\ref{inv:binding} completes the list. Each of these has an
executable analogue in at least one prototype, and the corresponding negative
tests are part of the recorded evidence.

\subsection{Adversarial failure modes and their defences}

The most dangerous implementation errors are not weak heuristics. They are
mistakes that quietly turn search evidence into a stronger assumption.

\begin{longtable}{@{}>{\raggedright\arraybackslash}p{0.30\textwidth}>{\raggedright\arraybackslash}p{0.64\textwidth}@{}}
\toprule
Failure mode & Required defence \\
\midrule
\endhead
Self-justifying lemma &
Prove candidates in a separate obligation; forbid dependency cycles unless a
joint induction principle explicitly justifies them. \\
Induction hypothesis applied to the wrong argument &
Keep the recursion index rigid; let only genuinely generalised binders be
instantiated, and record the permitted instantiation set as checked data. \\
Branch fact leakage &
Record split assumptions and scope; combine branches with a proof, never with
a global cache insertion. \\
Floating-point near-equality &
Recompute exact rational identities. A small residual is a rejection, not a
proof. \\
Numerical tolerance exploited in a weight &
Require the weight to be an exact rational \emph{of the right type} before
comparison; a weight of $(10^{20}+1)/10^{20}$ must not pass as $1$. \\
Wrong root problem &
Bind the certificate to the original reified expression and assumptions, never
to an oracle-provided replacement (Invariant~\ref{inv:binding}). \\
Missing subdivision region &
Derive child boxes from a full split tree and check both children; reject a
split point at an interval endpoint. \\
Duplicate keys in a decoded certificate &
Reject duplicate JSON object keys at parse time; two parsers disagreeing about
which value wins is a real differential hazard. \\
Vacuous replay success &
Fail when a required certificate family is absent from the results directory,
and mutate each decoded certificate to confirm the checker actually rejects
it. \\
Erased dependent type or cast &
Preserve typed denotation proofs and explicit transport obligations. \\
Numerical field proof used for integer division &
Use a source-semantics bridge, or reject the fragment. \\
Opaque model element used as a witness &
Reify an actual source term with the correct scope; otherwise reject the
proposal. \\
Incomplete external reconstruction &
Keep every hole as an unsolved obligation. No success until all are checked. \\
Stale state after metavariable assignment &
Invalidate affected caches, or restore a consistent snapshot containing the
old identities. \\
\bottomrule
\end{longtable}

\subsection{The proof-cost budget}

Search can produce a valid but impractically large certificate, and proof size
is a resource distinct from search time. Track certificate term count,
polynomial support size, rewrite-trace length, proof-DAG size, final proof-term
size, replay time, and axiom-audit time.

Share repeated polynomial normalisations and repeated nonnegativity proofs.
Emit a proof DAG topologically as local lemmas rather than re-expanding shared
subterms. Prefer sparse certificates and small rational coefficients when
several valid certificates are available --- but never alter a certificate
approximately in order to reduce its size, and never delete a dependency
because it looks unused in a heuristic trace while remaining inside the actual
proof term. Minimise a replay script only by rerunning and rechecking the
minimised version.

When a proof is rejected for size or replay budget, return a useful \unknown{}
diagnostic and preserve the candidate certificate as a debugging artefact. Do
not accept a partial proof, and do not silently raise the trust level in order
to finish the goal.

\subsection{The external certificate protocol}

Exchange structured data, never executable code. A problem request carries a
schema version, registered variable identifiers, exact rational coefficients,
normalised relation kinds, input bounds, and a resource envelope. A returned
certificate may reference only registered expressions and hypotheses; the Lean
side reconstructs the original semantics from its own registry.

The parser must reject malformed dimensions, duplicate or non-canonical
monomials where canonical form is required, negative exponents, oversized
integers, unknown references, invalid tree splits, and unsupported certificate
variants. It must not use \code{eval}, run a returned script, load arbitrary
shared libraries, or trust a theorem name supplied by the external solver.
Limits should be enforced \emph{before} expensive polynomial expansion as well
as after it, and search processes should be isolated from the proof-checking
process.

The packaged decoders demonstrate bounded rational parsing with numerator and
denominator bit-length caps, arity, term-count, total-degree, tree-depth and
file-size limits, and duplicate-key rejection. They are research tooling, not
hardened services for hostile uploads. A denial of service is not a soundness
failure, but it still makes a tactic unusable.
